% A139: smaller actual detector for the A134/A137 completed pair.
\subsection{A metabelian detector of order \(4096\) and class eight}

\begin{statement}{A.139}{Proved}
{The cyclic harmonic quotient already excludes the completed pair}
Retain the completed pair \((A,B)\) and ordinary centres \(x_p\)
of A.137, constructed from A.134.  Put
\[
 \begin{gathered}
 a=a_0a_1a_h=\frac{318039202706815117780707}{27871484},\qquad
 c=-\frac{338}{191},\qquad d=18512174429,\\
 \rho=\frac{34}{1088951437},\qquad
 F(T)=T(T-1)(T-\tfrac12).
 \end{gathered}
\tag{A.139.1}
\]
The connected degree-eight finite etale cover
\[
 D^\circ:\qquad W^8=\frac{t^2(1-t)}{64a^2}
                  \longrightarrow U
\tag{A.139.2}
\]
has geometric subgroup
\[
 H_D=\ker\bigl(\Pi\xrightarrow{\,2\kappa(t)+\kappa(1-t)\,}
                         \mathbf Z/8\bigr),\qquad
 K_D=\overline{H_D^2[H_D,H_D]}.
\tag{A.139.3}
\]
The actual finite geometric quotient is
\[
 Q_D=\Pi/K_D
       \simeq (\mathbf Z/2)^8\rtimes\mathbf Z/16,\qquad |Q_D|=4096,
\tag{A.139.4}
\]
where a generator of \(\mathbf Z/16\) cyclically permutes the eight
coordinates.  Its nilpotency class is exactly eight.

There is no continuous section of \(E_U/K_D\to G_\Q\) whose combined
Kummer class satisfies
\[
                 [\,t^2(1-t)\,]_8=[64a^2],
\tag{A.139.5}
\]
whose localization at every finite prime is geometrically conjugate
to the image of \(x_p\), and whose real localization is the image of an
ordinary or cuspidal real section.  The character in (A.139.5) is
defined on this quotient.
Consequently the actual arithmetic maps
\[
       E_U/\Pi''\longrightarrow E_U/K_D,\qquad
       E_U/\Gamma_9\Pi\longrightarrow E_U/K_D
\tag{A.139.6}
\]
exclude locally matching sections of the full metabelian quotient
and of the lower-central quotient of class eight for this completed
pair.  In particular the pair has no full Selmer section.
\end{statement}

\begin{proof}
\emph{The maps of curves and the proper Brauer class.}
The product-root map from A.134 is the literal morphism
\[
 \begin{gathered}
 C^\circ\longrightarrow D^\circ,\qquad
       (T,X_0,X_1,X_h)\longmapsto(T,W=X_0X_1X_h),\\
                         W^4=F(T)/a .
 \end{gathered}
\tag{A.139.7}
\]
Here the second model of \(D^\circ\) removes
\(T=0,1,\tfrac12,\infty\).
The isomorphism with (A.139.2), including its inverse, is
\[
     t=4T(1-T),\qquad
     W^8=\frac{t^2(1-t)}{64a^2},\qquad
     T=\frac12-\frac{4aW^4}{t}.
\tag{A.139.8}
\]
Indeed \(T(T-1)=-t/4\) and
\((T-\tfrac12)^2=(1-t)/4\), which prove both displayed
equations and both compositions.
The Kummer exponent at \(t=1\) is one, so the cover is geometrically
connected.  Its geometric group is cyclic of order eight.
In the degree-four model the three finite zeros of \(F\) and its
pole of order three are all fully ramified.  Riemann--Hurwitz gives
\(2g(D)-2=-8+4\cdot3=4\), hence \(g(D)=3\);
there are exactly four boundary points.

The class
\[
                        \alpha_D=(T-c,d)
\tag{A.139.9}
\]
extends to the smooth proper curve \(\overline D\).
This can be checked directly, without transferring properness
through (A.139.7).
At a point above \(T=c\) let \(w\) be its residue root.
The exact constants satisfy
\[
 \frac{F(c)}a=-\rho^2,\qquad
 \iota=\frac{w^2}{\rho},\qquad \iota^2=-1,\qquad
       ((1-\iota)w)^2=2\rho,\qquad
 \frac{d}{2\rho}=\left(\frac{1088951437}{2}\right)^2 .
\tag{A.139.10}
\]
Thus the residue field contains \(\sqrt d\) through the
displayed map.  The residue of (A.139.9) is zero there.
At infinity \(T\) has pole order four, so its quadratic residue
is again zero.  At every other point the first entry is a unit.
The unramified Brauer exact sequence therefore gives the claimed
proper extension.  Its pullback to \(C^\circ\) is precisely the
class of A.134.

\emph{All possible local sheets of this smaller cover.}
For each finite prime, the prescribed \(T_p\) gives a point of
\(D^\circ(\Q_p)\) by the product-root map.
For any point of \(D^\circ(\Q_p)\) lying above \(x_p\), the quadratic
identity in (A.139.8) gives exactly
\[
                              T=T_p\quad\hbox{or}\quad T=1-T_p.
\tag{A.139.11}
\]
Since \(F(1-T_p)=-F(T_p)\), the second possibility exists exactly
when \(-1\) is a fourth power in \(\Q_p\).  This statement follows
by taking the ratio of the two nonzero fourth-root coordinates;
it is independent of the choice of the first root.
It is impossible at \(2\), and at an odd prime occurs exactly for
\(p\equiv1\pmod8\).

Let
\[
                  S=\{2,3,7,13,17,23,139,191,1119169\}.
\tag{A.139.12}
\]
At every odd \(p\notin S\), the evaluation of \(\alpha_D\) is zero
on every ordinary \(D(\Q_p)\)-point, not just on the prescribed
fibre.  All \(a,\rho,d,c\) and old-puncture differences are
\(p\)-integral units.  If \(v_p(T)<0\), then
\(4v_p(W)=3v_p(T)\), so \(4\mid v_p(T)\);
also \(v_p(T-c)=v_p(T)\).
The unramified quadratic character of the unit \(d\) therefore
evaluates to zero.  If \(T\) is integral and \(T-c\) is a unit,
both entries are units and give zero.  In the remaining case,
reducing the fourth-root equation modulo \(p\) gives a nonzero
root of \(Z^4=-\rho^2\).  Formula (A.139.10), now in the residue
field, makes \(d\) a square modulo \(p\); Hensel's lemma splits
the local quadratic character.  This also covers the actual
points with \(T=c\), by properness and continuity.

Inside \(S\), the only odd primes allowing the second sheet
in (A.139.11) are \(17\) and \(r=1119169\).
At \(17\), A.134 fixes
\[
 T_{17}=1+a_1\,17^4,\qquad c\equiv9\pmod{17},\qquad
       T_{17}-c\equiv9,\quad 1-T_{17}-c\equiv8\pmod{17}.
\tag{A.139.13}
\]
Both residues are squares: \(9=3^2\) and \(8=5^2\) modulo \(17\).
At \(r\), its fixed centre is \(T_r=\tfrac12+a_h\) with
\(v_r(a_h)=2\).  Both first entries have the same nonzero residue
\[
            \frac12-c=\frac{867}{382}\pmod r,
\tag{A.139.14}
\]
which is a square by A.134's executed residue calculation
(equivalently its squareclass is \(1146\), and \((1146/r)=1\)).
Hence the flipped sheet also has zero invariant at these two
primes.
At every other prime in \(S\) only \(T_p\) occurs, and A.134 gives
the evaluation through (A.139.7).
We have proved, for every possible local sheet above \(x_p\),
\[
       \operatorname{inv}_p\alpha_D=
       \begin{cases}
          1/2,&p=2,\\
          0,&p\ne2.
       \end{cases}
\tag{A.139.15}
\]
At the real place \(d>0\), so the class is split at every ordinary
or proper boundary point.

\emph{The actual finite quotient and its exact class.}
Use the standard free meridians \(x,y\) around \(0,1\).
The character in (A.139.3) takes them to \(2,1\) modulo eight.
The change of free generators
\[
                         u=xy^{-2},\qquad v=y
                         \quad(x=uv^2,\ y=v)
\tag{A.139.16}
\]
is an explicit Nielsen isomorphism.
The Schreier basis of \(H_D\) is
\[
        u_j=v^juv^{-j}\quad(0\le j<8),\qquad z=v^8.
\tag{A.139.17}
\]
Consequently \(H_D/K_D\) is the elementary abelian group on these
nine independent classes.
Conjugation by \(v\) sends \(u_j\) to \(u_{j+1}\), with indices
modulo eight: at the wraparound, \(v^8=z\) commutes with all of
\(H_D/K_D\).
Also \(v^{16}=z^2=1\).
These relations give the map to the semidirect product in
(A.139.4), \(x\mapsto u_0v^2,\ y\mapsto v\).
Conversely that semidirect product has independent commuting
classes \(u_0,\ldots,u_7,v^8\), all of order two.
The image of \(H_D\) is therefore exactly its mod-two abelianization.
The two maps prove that the kernel is precisely \(K_D\), and give
the order \(2^9\cdot8=4096\).
The geometric kernel of the Kummer character is stable under the
arithmetic action, which multiplies that character by the
cyclotomic unit.  Since \(K_D\) is characteristic in \(H_D\),
it is normal in \(E_U\).

For the exact class, regard the lamp group as
\[
 L=\F_2[C_8]=\F_2[\tau]/(\tau^8),\qquad \tau=v-1.
\tag{A.139.18}
\]
Conjugation by \(v\) is multiplication by \(1+\tau\).
The commutator subgroup of \(Q_D\) is \(\tau L\):
the lamp group and the cyclic quotient are abelian, and
\([u_0,v]=(1+v)u_0\) generates exactly that submodule under
conjugation.  The actual recurrence is
\[
       \Gamma_kQ_D=\tau^{k-1}L\quad(2\le k\le8),\qquad
       \Gamma_9Q_D=0.
\tag{A.139.19}
\]
Thus the dimensions are \(7,6,5,4,3,2,1,0\).
The left-normed commutator obtained from \([x,y]\) by six further
commutators with \(y\) equals
\(\tau^7u_0=\sum_{j=0}^7u_j\ne0\).
This proves exact class eight without enumeration.
Since \(\Pi'\subset H_D\),
\(\Pi''\subset[H_D,H_D]\subset K_D\).
Together with the computed class this gives the two actual maps
in (A.139.6).
Finally the integral combined Kummer character takes \(H_D\) into
\(8\widehat\Z(1)\) and \(K_D\) into \(16\widehat\Z(1)\).
In particular (A.139.5) is a character of \(E_U/K_D\).

\emph{Finite descent and reciprocity.}
Suppose a section with the asserted data is given.
Condition (A.139.5) says exactly that its pullback of the root
torsor (A.139.2) is trivial.
The finite-cover correspondence supplies a fixed geometric sheet,
and after geometric conjugation the section lands in
\[
              E_D/K_D\subset E_U/K_D,\qquad
              E_D=\pi_1(D^\circ).
\tag{A.139.20}
\]
Choose a lift of \(\alpha_D\) to
\(H^2_{\rm et}(\overline D,\mu_2)\), restrict it to \(D^\circ\),
and use the exact group/etale comparison in (A.137.6).
The resulting central extension of \(E_D\) splits over the free
geometric group \(H_D\).
For a splitting \(j\), its conjugation discrepancy is a character
\(H_D\to\F_2\), which kills \(K_D\).
The actual normal subgroup \(j(K_D)\) and fibre-product isomorphism
of (A.137.9), with \(H_D,K_D,E_D\) in place of \(H,K,E\),
therefore descend this class to
\[
                               e_D\in H^2(E_D/K_D,\mu_2).
\tag{A.139.21}
\]
This construction uses no rational point of \(D^\circ\) and no
arithmetic splitting of \(E_D\).

Pulling (A.139.21) back along (A.139.20) gives a rational Brauer
class.  At each finite prime, the local fixed sheet of the actual
cover is a point of \(D^\circ(\Q_p)\) above \(x_p\).
The comparison and fibre-product maps identify its local value
with \(\alpha_D\) at that point.  All possible sheets were covered
in (A.139.15).  At the real place, an ordinary image gives ordinary
evaluation; for a cuspidal image, the inertia orbit is a real point
of the proper normalization.  Since the chosen \(\mu_2\) class
comes from \(\overline D\), its tangential pullback is that proper
evaluation, exactly as in A.137.  It is zero.
The sum of local invariants is therefore \(1/2\), contradicting
reciprocity for this actual class in \(\operatorname{Br}(\Q)[2]\).
This proves the finite-quotient exclusion.

The completed pair from A.137 satisfies
\[
             A^2B=64(A_0A_1A_h)^2
                    =[64(a_0a_1a_h)^2]\pmod{(\Q^\times)^8},
\tag{A.139.22}
\]
because each old completed coordinate has its indicated mod-four
class.  Its full local Kummer coordinates force the ordinary
centres \(x_p\), and exclude every local cusp, by A.137.
Projection of a full Selmer section, or a locally matching section
of either quotient in (A.139.6), would give the excluded section.
This proves all assertions.
\end{proof}

The retained certificate
\path{code/harmonic_cyclic_detector_certificate.py}, its output,
and \path{code/harmonic_cyclic_detector_data.json}
check the exact constants and both curve maps, all \(611\) retained
local centres and every available flipped sheet, and the actual
eight-vertex path-chain model of \(Q_D\).
The latter has \(16\) directed edges, cycle rank nine, exactly
\(4096\) elements, and lower-central dimensions
\(7,6,5,4,3,2,1,0\).  Its specific nonzero weight-eight commutator
is the sum of all eight \(x\)-edges.
